%% notes-tex/010-additive-baselines/sections/4-models.tex

\section{The models\secstatus{todo}}
\label{sec:models}

\subsection{Where each baseline comes from\secstatus{todo}}
\label{sec:provenance}

Two of the six baselines are taken from the Visani preprint. The rest were built
here for a confounder that the protein setting does not have. Naming which is which matters, because a baseline borrowed from a published
critique carries that critique's authority and one invented here does not.

%% SOURCE: structural; the two external rows cite the preprint by DOI.
\begin{table}[htbp]
\centering
\begin{tabular}{lll}
\hline
Model & Origin & Corresponds to \\
\hline
B0 & standard & intercept-only null \\
B1 & Visani et al.\ 2026 & their ridge-regularized additive model \\
B2 & this document & B1 extended to observed pairs \\
B3 & this document & unregularized counterpart of B4 \\
B4 & this document & the Kuzmin screen-structure confounder \\
B5 & Visani et al.\ 2026 & their zero-hidden-layer control, run in reverse \\
\hline
\end{tabular}
\caption{Provenance of the six baselines.}
\label{tab:modelorigin}
\end{table}

The two external rows come from Visani, Verma and DeWitt, \emph{Additive
baselines furnish no evidence for epistasis learning by MULTI-evolve}, bioRxiv
\texttt{2026.04.23.719915}
(\url{https://doi.org/10.1101/2026.04.23.719915}), which is itself a reply to
Tran et al., \emph{Science} 2026, \texttt{10.1126/science.aea1820}
(\url{https://doi.org/10.1126/science.aea1820}). B1 is their additive null
transplanted onto gene deletions. B5 inverts their zero-hidden-layer control:
where they removed nonlinearity from the published architecture, B5 adds the
smallest amount of nonlinearity to the additive feature space, so the two
bracket the same gap from opposite sides.

B3 and B4 have no counterpart in the preprint because the confounder they measure
is specific to a synthetic genetic array. A protein engineering campaign has no
analog of a query double mutant recurring across thousands of measured
variants. B4 turns out to carry the result that matters most here, so its being
ours rather than theirs is worth stating plainly.

\subsection{Notation and the interaction operator\secstatus{todo}}

Let $G$ be the 4{,}352 genes appearing in the build and let record $n$ have
perturbed set $S_n \subset G$ with $|S_n| = 3$. Encode it as
$x_n \in \{0,1\}^{|G|}$ with $x_{n,g} = 1$ exactly when $g \in S_n$, so that
$\lVert x_n \rVert_1 = 3$ for every record. Write $y_n$ for the trigenic
interaction score of Eq.~\ref{eq:tau}.

A model is a set function $f : 2^{G} \to \mathbb{R}$. Its capacity to represent
interaction is measured by finite differences. For distinct genes $g, h, k$ and
any $S$ disjoint from them, define
%
\begin{align}
  \Delta_g f(S) &= f(S \cup \{g\}) - f(S), \\
  \Delta_{gh} f(S) &= f(S \cup \{g,h\}) - f(S \cup \{g\}) - f(S \cup \{h\}) + f(S),
  \label{eq:delta2}\\
  \Delta_{ghk} f(S) &= \sum_{U \subseteq \{g,h,k\}} (-1)^{3-|U|} f(S \cup U).
  \label{eq:delta3}
\end{align}
%
$\Delta_{gh} f \equiv 0$ says the model cannot express any pairwise interaction;
$\Delta_{ghk} f \equiv 0$ says it cannot express any three-way interaction. These
are the properties that make a model a null, and they are structural, holding for
every parameter setting rather than for a fitted one.

\subsection{B0, the train mean\secstatus{todo}}

$\hat{y} = \bar{y}_{\mathrm{train}}$. Every finite difference vanishes. Its mean
squared error is the label variance, and it fixes the zero point against which
all reductions are read.

\subsection{B1, the additive per-gene ridge\secstatus{todo}}

%
\begin{equation}
  \hat{y}(S) \;=\; \beta_0 + \sum_{g \in S} \beta_g ,
  \qquad
  \hat{\beta} \;=\; \arg\min_{\beta} \sum_{n \in \mathrm{train}}
  \bigl(y_n - \beta_0 - x_n^{\top}\beta\bigr)^2 + \alpha \lVert \beta \rVert_2^2 ,
  \label{eq:b1}
\end{equation}
%
with $\beta_0$ unpenalized and set to $\bar{y}_{\mathrm{train}}$. Writing $X$
for the $n_{\mathrm{train}} \times |G|$ matrix whose rows are the $x_n$ and
$\tilde{y} = y - \bar{y}_{\mathrm{train}}$ for the centered labels, the
minimizer is the ridge normal equation
%
\begin{equation}
  \hat{\beta} \;=\; \bigl(X^{\top} X + \alpha I\bigr)^{-1} X^{\top} \tilde{y} ,
  \label{eq:b1closed}
\end{equation}
%
which the script evaluates by LSQR on the augmented system
$[X;\, \sqrt{\alpha}\, I]\,\beta = [\tilde{y};\, 0]$ rather than by forming the
inverse. Substituting Eq.~\ref{eq:b1} into Eq.~\ref{eq:delta2} gives
$\Delta_{gh}\hat{y} = 0$ and hence $\Delta_{ghk}\hat{y} = 0$ identically. B1 is
the direct analog of the additive null in the Visani preprint.

The design is well conditioned despite having 4{,}352 coefficients: with three
ones per row and 301{,}386 training rows there are 904{,}158 gene observations,
a median of 9 per gene in the training split with a maximum of 3{,}430. The
ridge penalty is chosen on validation over
$\alpha \in \{10^{-2}, 10^{-1}, 1, 3, 10, 30, 100, 300, 1000\}$ and the selected
value is then reported on test.

\subsection{B2, additive plus recurring gene pairs\secstatus{todo}}

Let $P$ be the set of unordered gene pairs occurring at least five times in the
training split, $|P| = 420$. Adding one coefficient per such pair,
%
\begin{equation}
  \hat{y}(S) \;=\; \beta_0 + \sum_{g \in S} \beta_g
  + \sum_{\{g,h\} \subseteq S,\; \{g,h\} \in P} \gamma_{gh} ,
  \label{eq:b2}
\end{equation}
%
gives $\Delta_{gh}\hat{y} = \gamma_{gh}$ for a pair in $P$ and zero otherwise.
B2 therefore has second-order capacity on observed pairs. It still has
\emph{no} third-order capacity: $\Delta_{ghk}\hat{y} = 0$ identically. Since the
label is itself a third-order residual, B2 remains a null with respect to the
quantity being predicted.

The fit is the same ridge as B1 on the wider design $X = [X_G \;\; X_P]$, where
$X_P$ is the $n_{\mathrm{train}} \times |P|$ pair-incidence matrix, so
$(\hat{\beta}, \hat{\gamma})$ is Eq.~\ref{eq:b1closed} with $X$ replaced by
$[X_G \;\; X_P]$ and one penalty $\alpha$ on both blocks. The pair block has a
closed form of its own that says what B2 learns. Every record contains exactly
one recurring pair (Section~\ref{sec:disjoint}), so the stationarity condition
of the ridge objective in $\gamma_{gh}$ involves only that pair's records, and
%
\begin{equation}
  \hat{\gamma}_{gh} \;=\; \frac{\sum_{n \in \mathrm{train},\; \{g,h\} \subseteq S_n}
  \bigl(y_n - \beta_0 - x_n^{\top}\hat{\beta}\bigr)}{n_{gh} + \alpha} ,
  \label{eq:b2gamma}
\end{equation}
%
with $n_{gh}$ the number of training triples carrying the pair. Each pair
coefficient is the mean residual of the additive model over that query double's
screen, shrunk toward zero by $\alpha / (n_{gh} + \alpha)$. With $n_{gh}$ in the
hundreds to thousands and the selected $\alpha$ in the units, the shrinkage is
at most a few percent, so B2 is B1 plus one nearly unshrunk offset per screen.

\subsection{B3, the empirical screen mean\secstatus{todo}}
\label{sec:b3}

B3 has no fitted parameters. Its per-gene and per-pair quantities are means of
the training label over the records that contain the gene or the pair,
%
\begin{equation}
  \bar{y}_g \;=\; \frac{\sum_{n \in \mathrm{train}} x_{n,g}\, y_n}{\sum_{n \in \mathrm{train}} x_{n,g}},
  \qquad
  \bar{y}_{gh} \;=\; \frac{\sum_{n \in \mathrm{train}} x_{n,g}\, x_{n,h}\, y_n}{\sum_{n \in \mathrm{train}} x_{n,g}\, x_{n,h}} ,
  \label{eq:b3means}
\end{equation}
%
so $\bar{y}_g$ averages $y$ over exactly the training triples in which gene $g$
is one of the three deletions, and $\bar{y}_{gh}$ over those carrying both $g$
and $h$. The prediction backs off through three levels,
%
\begin{equation}
  \hat{y}(S) \;=\;
  \begin{cases}
    \operatorname{mean}\{\bar{y}_{gh} : \{g,h\} \subseteq S,\; n_{gh} \ge 5\}
      & \text{if any pair in } S \text{ recurs,} \\
    \operatorname{mean}\{\bar{y}_{g} : g \in S,\; n_g \ge 5\}
      & \text{otherwise, if any gene in } S \text{ has 5 training records,} \\
    \bar{y}_{\mathrm{train}} & \text{otherwise,}
  \end{cases}
  \label{eq:b3}
\end{equation}
%
with $n_g$ and $n_{gh}$ the training counts. On the random split every test
record carries a recurring pair, so the first case always fires and B3 is the
per-screen mean of the label. On a query-pair-disjoint split no test record
carries a recurring pair, the first case never fires, and B3 becomes the mean of
its three gene means.

\subsection{B4, the screen indicator\secstatus{todo}}
\label{sec:b4}

B4 drops the per-gene terms from Eq.~\ref{eq:b2} and keeps only the pair terms,
%
\begin{equation}
  \hat{y}(S) \;=\; \beta_0 + \sum_{\{g,h\} \subseteq S,\; \{g,h\} \in P} \gamma_{gh} ,
  \qquad
  \hat{\gamma}_{gh} \;=\; \frac{\sum_{n \in \mathrm{train},\; \{g,h\} \subseteq S_n}
  \bigl(y_n - \bar{y}_{\mathrm{train}}\bigr)}{n_{gh} + \alpha} ,
  \label{eq:b4}
\end{equation}
%
where the second expression is Eq.~\ref{eq:b2gamma} with no gene block. Because
the recurring pairs are the Kuzmin query doubles, B4 is a per-screen offset: for
a triple made of a query pair and an array gene, the array gene contributes
nothing at all. Whatever B4 achieves is therefore an upper bound on how much of
the metric is explained by knowing which screen a record came from.

Eq.~\ref{eq:b4} at $\alpha \to 0$ is $\bar{y}_{gh}$ of Eq.~\ref{eq:b3means},
so on the random split, where the first case of Eq.~\ref{eq:b3} always fires,
B3 and B4 are the same predictor up to the shrinkage $\alpha/(n_{gh}+\alpha)$,
which Pearson does not see. Table~\ref{tab:heldout} reports them equal to three
decimals for that reason. On a disjoint split they part company: B4 is
structurally zero and B3 falls back to gene means.

\subsection{B5, nonlinearity on the same features\secstatus{todo}}

Give each gene a free vector $e_g \in \mathbb{R}^{128}$, sum over the perturbed
set and read out with a two-hidden-layer network $\phi$,
%
\begin{equation}
  \hat{y}(S) \;=\; \phi\!\left(\sum_{g \in S} e_g\right),
  \qquad
  \phi : \mathbb{R}^{128} \to \mathbb{R}.
  \label{eq:b5}
\end{equation}
%
This is permutation invariant and uses precisely B1's feature space, the identity
of the three genes and nothing more. Because $\phi$ is nonlinear,
$\Delta_{gh}\hat{y}$ and $\Delta_{ghk}\hat{y}$ are not identically zero, so B5 is
not a null. It is the smallest departure from B1 that has interaction capacity,
and it isolates what capacity alone buys.

Comparing Eq.~\ref{eq:b5} with Eq.~\ref{eq:cgt} shows they are the same
construction. Both give every gene a free vector, pool over the three perturbed
genes in a permutation-invariant way, and pass the pooled vector through a small
network. The transformer's pooling is richer, since the context $c_S$ is
produced by attention over the perturbed set rather than by a plain sum, and its
per-gene vectors are $Z = T(E)$ rather than a raw table. Trained with three
seeds, B5 reports a spread rather than a single number.

\subsection{Summary of capacity\secstatus{todo}}

%% SOURCE: structural, derived above from the model equations; no data involved
\begin{table}[htbp]
\centering
\footnotesize
\begin{tabular}{lp{0.22\textwidth}p{0.24\textwidth}lcc}
\hline
Model & Per-gene quantity & Set aggregation & Fit by & $\Delta_{gh} \not\equiv 0$ & $\Delta_{ghk} \not\equiv 0$ \\
\hline
B0 & none & none & train mean & no & no \\
B1 & free scalar $\beta_g$ & sum & ridge, Eq.~\ref{eq:b1closed} & no & no \\
B2 & free scalar $\beta_g$ & sum, plus $\gamma_{gh}$ per recurring pair & ridge, Eq.~\ref{eq:b2gamma} & yes & no \\
B3, pair case & none & mean of $\bar{y}_{gh}$ over the recurring pairs in $S$ & label means, Eq.~\ref{eq:b3means} & yes & no \\
B3, gene case & $\bar{y}_g$, the label mean over training triples containing $g$ & mean over $S$ & label means, Eq.~\ref{eq:b3means} & no & no \\
B4 & none & $\gamma_{gh}$ per recurring pair & ridge, Eq.~\ref{eq:b4} & yes & no \\
B5 & free 128-vector $e_g$ & sum, then MLP & gradient descent & yes & yes \\
CGT & free 180-vector via $T(E)$ & attention over $S$, mean, MLP & gradient descent & yes & yes \\
\hline
\end{tabular}
\caption{What each model can represent. B3 is listed by the case of
Eq.~\ref{eq:b3} that fires: the pair case on the random split, the gene case on
a query-pair-disjoint split. The last two columns are structural properties of
the hypothesis class, not fitted results. Only B5 and the cell graph transformer
can express a three-way interaction, which is what the label is.}
\label{tab:capacity}
\end{table}
